\chapter{Patient Assessment}\label{ch:assessment}

Every first aid response should follow a logical assessment process. Do not rush to treat an injury before you know whether a life-threatening condition exists. The assessment framework has two parts: a \textbf{primary survey} (find and fix life threats) and a \textbf{secondary survey} (identify all other injuries and gather history).

\section{The Primary Survey}

The primary survey answers one question: \emph{Is anything immediately life-threatening?} Work through it in order, and stop to treat anything you find before moving on.

\begin{enumerate}[leftmargin=*]
\item \textbf{Scene safety} -- Confirm the scene is safe before approaching. Do not become another victim.
\item \textbf{Responsiveness} -- Tap the shoulders and shout, ``Are you OK?'' A person who opens their eyes or responds is responsive.
\item \textbf{Activate EMS} -- If unresponsive, call emergency services immediately (or send someone to call while you stay).
\item \textbf{Airway \& Breathing} -- Look, listen, and feel for normal breathing for 5--10 seconds. If the person is not breathing normally (or only gasping), begin CPR.
\item \textbf{Circulation} -- Check for severe bleeding and control it with direct pressure. If no pulse and not breathing, begin chest compressions.
\end{enumerate}

\textbf{WARNING:} The mnemonic \textbf{CAB} (Compressions, Airway, Breathing) reminds you that in cardiac arrest, chest compressions come first: give 30 compressions before the two rescue breaths.

\quicknote{\textbf{One flow, all emergencies:} Unresponsive and not breathing normally $\rightarrow$ CPR. Unresponsive but breathing normally $\rightarrow$ recovery position and monitor. Responsive $\rightarrow$ secondary survey, treat injuries, and monitor.}

\section{The Secondary Survey}

Perform the secondary survey on any person who is responsive, or on an unresponsive person whose life threats have been stabilized while you wait for help. Its purpose is to find injuries and conditions you have not yet noticed.

\subsection{Head-to-Toe Examination}

Work systematically from head to feet, comparing one side of the body with the other:

\begin{itemize}[leftmargin=*]
\item \textbf{Head}: scalp lumps, depressions, bleeding, bruising around eyes or behind ears
\item \textbf{Neck}: pain, tenderness, swelling
\item \textbf{Chest}: pain on breathing, visible deformities, equal chest rise
\item \textbf{Abdomen}: pain, rigidity, distension
\item \textbf{Pelvis}: pain or instability (do not rock the pelvis to test it)
\item \textbf{Arms and legs}: pain, deformity, swelling, and \textbf{nerve check} (sensation, movement, and circulation -- pulse and capillary refill) distal to any injury
\item \textbf{Back}: check last, only if the person can be rolled safely (log roll) without spinal risk
\end{itemize}

Use the memory aid \textbf{DCAP-BTLS} to remember what to look for at each location: \emph{Deformities, Contusions (bruises), Abrasions, Punctures/penetrations, Burns, Tenderness, Lacerations, Swelling}.

\subsection{SAMPLE History}

Once you know there are no life threats, gather a focused history. SAMPLE is the standard memory aid:

\begin{description}[leftmargin=!]
\item[S -- Symptoms:] What is the main complaint? Describe the symptoms.
\item[A -- Allergies:] Any known allergies (foods, medications, insects, latex)?
\item[M -- Medications:] What medications does the person take, including prescription, over-the-counter, and supplements?
\item[P -- Past medical history:] Known conditions (diabetes, heart disease, epilepsy, asthma)?
\item[L -- Last meal:] What and when did they last eat or drink?
\item[E -- Events:] What happened leading up to the emergency?
\end{description}

\subsection{OPQRST for Pain or Discomfort}

To characterize a symptom such as chest or abdominal pain, use OPQRST:

\begin{description}[leftmargin=!]
\item[O -- Onset:] Did the pain start suddenly or gradually?
\item[P -- Provocation:] What makes it better or worse (movement, breathing, rest)?
\item[Q -- Quality:] Sharp, dull, crushing, burning, pressure?
\item[R -- Radiation:] Does it spread to the arm, back, jaw, or elsewhere?
\item[S -- Severity:] Rate it on a 0--10 scale.
\item[T -- Time:] How long has it been present?
\end{description}

\section{Vital Signs}

Vital signs help you recognize how sick or injured a person is and detect deterioration. Record your findings and the time; report them to responders.

\subsection{Level of Consciousness -- AVPU}

\begin{description}[leftmargin=!]
\item[A -- Alert:] Responsive and oriented.
\item[V -- Voice:] Responds when you speak loudly or shout.
\item[P -- Pain:] Responds only to a painful stimulus (e.g., a firm pinch to the shoulder).
\item[U -- Unresponsive:] No response to voice or pain.
\end{description}

A decreasing level of consciousness is a red flag -- reassess frequently and call emergency services if not already done.

\subsubsection{Glasgow Coma Scale (GCS)}
The Glasgow Coma Scale provides a standardized way to assess level of consciousness in head injury patients. Score each component and add the totals (range 3--15):

\begin{center}
\small
\begin{tabularx}{\textwidth}{@{}>{\raggedright\arraybackslash}lX>{\raggedright\arraybackslash}r@{}}
\hline
\textbf{Component} & \textbf{Response} & \textbf{Score} \\
\hline
\textbf{Eye Opening} & Spontaneous & 4 \\
& To speech & 3 \\
& To pain & 2 \\
& None & 1 \\
\textbf{Verbal Response} & Oriented & 5 \\
& Confused & 4 \\
& Inappropriate words & 3 \\
& Incomprehensible sounds & 2 \\
& None & 1 \\
\textbf{Motor Response} & Obeys commands & 6 \\
& Localizes pain & 5 \\
& Withdraws from pain & 4 \\
& Abnormal flexion (decorticate) & 3 \\
& Extension (decerebrate) & 2 \\
& None & 1 \\
\hline
\end{tabularx}
\end{center}

Quick reference: GCS 13--15 = mild head injury, 9--12 = moderate, 3--8 = severe (coma). Any GCS below 13 after head trauma warrants urgent medical evaluation. Record the score and time; repeat every 15 minutes in unstable patients.

\subsection{Normal Ranges by Age}

\begin{center}
\small
\begin{tabular}{@{}llll@{}}
\hline
\textbf{Age Group} & \textbf{Heart Rate (bpm)} & \textbf{Respiratory Rate (per min)} & \textbf{Blood Pressure (systolic)} \\
\hline
Infant (0--1 y) & 100--160 & 30--60 & 70--100 \\
Child (1--8 y) & 80--120 & 20--30 & 80--110 \\
Child (8--12 y) & 70--110 & 16--25 & 90--120 \\
Adult & 60--100 & 12--20 & 90--120 \\
\hline
\end{tabular}
\end{center}

\quicknote{These ranges are reference values for adults and children at rest. A pulse that is very fast, very slow, or irregular, or breathing that is very fast, very slow, or labored, should raise concern. Trends matter more than single readings -- always record serial values.}

\subsection{Other Observations}

\begin{itemize}[leftmargin=*]
\item \textbf{Skin}: color (pale, flushed, blue/gray), temperature, and moisture (clammy, dry)
\item \textbf{Capillary refill}: press a fingernail bed for 2 seconds; color should return within 2 seconds
\item \textbf{Pupils}: equal size and reaction to light (unequal pupils after head injury is a red flag)
\item \textbf{Skin bleeding points}: check for injuries you may have missed
\end{itemize}

\section{Recognition and Management of Shock}

Shock is a state in which the body cannot deliver enough oxygen to the vital organs. It can follow severe bleeding, burns, infection, allergic reactions, heart problems, or spinal injury. A person in shock can deteriorate quickly and may die even from injuries that look survivable.

\subsection{Signs of Shock}

\begin{itemize}[leftmargin=*]
\item Pale, cool, clammy skin
\item Rapid, weak pulse
\item Rapid, shallow breathing
\item Thirst, nausea, or vomiting
\item Restlessness, anxiety, or confusion
\item Altered level of consciousness (drowsiness, disorientation)
\item Weakness and dizziness, especially when sitting or standing
\end{itemize}

\subsection{Management}

\begin{enumerate}[leftmargin=*]
\item Call emergency services immediately.
\item Treat the cause you can see: control bleeding, immobilize fractures, keep the person warm.
\item Lay the person flat on their back -- this is the recommended position for shock.
\item Raise the legs about 6--12 inches (15--30 cm) ONLY for non-traumatic causes (e.g., simple fainting, dehydration) and ONLY if the position does not cause pain or worsen symptoms. Do NOT raise the legs if trauma is suspected, if the person has a head injury, breathing difficulty, or heart problems, or if the position is uncomfortable.
\item Loosen tight clothing.
\item Cover with a blanket or coat to prevent heat loss, but do not overheat.
\item Do NOT give food or drink (the person may require surgery). If thirsty, moisten their lips only.
\item Monitor breathing and level of consciousness continuously. Be ready to begin CPR.
\end{enumerate}

\textbf{WARNING:} Do NOT raise the legs of a person with chest pain, breathing difficulty, a head injury, or suspected trauma. Keep them in the position of comfort. Never leave a person in shock alone.

\section{Continuous Monitoring and Documentation}

Emergencies change. Reassess every few minutes:

\begin{itemize}[leftmargin=*]
\item Recheck responsiveness, breathing, and pulse.
\item Repeat vital signs and note the time of each reading.
\item Watch for skin changes and worsening pain.
\item Record everything: time of collapse, events, medications given, and all observations. Report this information to emergency responders.
\end{itemize}

\subsection{Handover to Emergency Responders (SBAR)}\label{sec:sbar}

When professional help arrives, give a clear, brief report. The SBAR format covers the essentials:

\begin{itemize}[leftmargin=*]
  \item \textbf{S -- Situation}: What is happening now? ``Adult male, unresponsive, not breathing.''
  \item \textbf{B -- Background}: What happened? ``Found collapsed in the kitchen about 10 minutes ago. Has diabetes, no known allergies.''
  \item \textbf{A -- Assessment}: What did you find? ``No response to voice or pain. Breathing stopped, no pulse felt (if trained to check). Skin pale and cool.''
  \item \textbf{R -- Recommendation / Response}: What was done and what is needed? ``CPR started 8 minutes ago, 30:2, AED attached, one shock delivered. Needs immediate transport.''
\end{itemize}

Keep the report under 30 seconds and let the responders ask questions. Hand over any written notes, medications given, and the person's belongings. Do not leave the person unattended; stay until the responders take over and tell you exactly when they assume care.

\section{Review Questions}

\begin{enumerate}[leftmargin=*]
\item What are the steps of the primary survey, in order, and when should you interrupt them to treat?
\item A person is unresponsive but breathing normally. What should you do?
\item What does the SAMPLE history stand for, and why is it useful?
\item How does the AVPU scale help you track a person's condition over time?
\item List three signs of shock and the first three things you should do for a person in shock.
\end{enumerate}

\index{primary survey}
\index{secondary survey}
\index{head-to-toe examination}
\index{DCAP-BTLS}
\index{SAMPLE history}
\index{OPQRST}
\index{vital signs}
\index{AVPU}
\index{capillary refill}
\index{shock}
\index{patient assessment}
\index{SBAR handover}
\index{handover to EMS}
